%
% Conclusion.tex
%
% Copyright (C) 2025 by Inspirefly.
%
% 2.305 GHz Circular Patch Antenna Documentation
%

\chapter{Conclusion}

A 2.305~GHz circular microstrip patch antenna for the PocketQube parent board has been sized, simulated, and laid out in a reproducible flow that links analytical cavity theory, MATLAB Antenna Toolbox exploration, ANSYS HFSS full-wave verification, and KiCAD 10.0 fabrication data.

The cavity model (Eqs.~\ref{eq:fres}--\ref{eq:ae}) gives a physical radius of $19.87$~mm on the final FR-408HR stack ($\epsilon_r=3.68$, $h=0.1$~mm prepreg), within $0.5$~mm of the HFSS-tuned radius used in \texttt{CircularPatchV3.aedt}. The early rectangular $\lambda_g/2$ estimate ($L\approx 17.8$~mm on RO4350B) and the hexagonal variant were evaluated and superseded, with provenance retained on branch \texttt{to-be-reviewed\_V0.1}. The MATLAB particle-swarm search (\path{matlab/MPA_opt_swrm.m}) demonstrated the limits of a single-layer lossless model, converging to a low-impedance minimum ($Z\approx 0.3$~$\Omega$, RL${}\approx 0$~dB) that confirmed the need for direct HFSS tuning of patch radius and feed offset.

The HFSS model \texttt{CircularPatchV3.aedt} implements two orthogonal probe feeds for circular polarization. When driven in quadrature through the external APS-02-WBM-LP-DF-CC divider, simulation predicts $S_{11} < -10$~dB per port at 2.305~GHz, inter-port isolation below $-18$~dB, gain near 4--5~dBi, and axial ratio below 3~dB at boresight. Geometry was exported and converted to the KiCAD footprint \texttt{Inspirefly:CircularPatchAntenna}. The final board can be found in (\path{KiCAD/antennaBoard/antennaBoard.kicad_pcb}. All sources (ANSYS .aedt projects in \path{ansys/}, MATLAB scripts and logs in \path{matlab/}, KiCAD schematics, PCB, footprints, and DXF/STEP exports in \path{KiCAD/}, and this document in \path{doc/}) are versioned in \path{Pocketqube-parent-antennaBoard} git.

